\paragraph{Sơ đồ tuần tự Đánh dấu cuộc trò chuyện (Bookmark Flow)}

Luồng nghiệp vụ này mô tả quy trình người dùng thực hiện đánh dấu (bookmark) một tin nhắn hoặc một cuộc trò chuyện cụ thể, sắp xếp theo thư mục tự chọn và đính kèm ghi chú cá nhân vào Neon Database.

\begin{figure}[H]
\centering
\includegraphics[scale = 0.18]{pictures/UML-Sequence-21.png}
\caption{Sơ đồ tuần tự Đánh dấu cuộc trò chuyện}
\label{fig:seq_bookmark}
\end{figure}

\subparagraph{Các thành phần tham gia}

\begin{table}[H]
\centering
\renewcommand{\arraystretch}{1.3}
\begin{tabular}{|l|l|p{5cm}|}
\hline
\textbf{Thành phần} & \textbf{Phân loại} & \textbf{Vai trò trong hệ thống} \\ \hline
User & Actor & Người dùng chọn tin nhắn cần lưu trữ và cấu hình thư mục. \\ \hline
Web Frontend & Component & Giao diện Next.js hiển thị form nhập ghi chú, chọn thư mục lưu trữ và hiển thị kết quả. \\ \hline
API Gateway (Kong) & Component & Định tuyến yêu cầu lưu bookmark qua phương thức POST. \\ \hline
Conversation Service & Microservice & Tiếp nhận yêu cầu, kiểm tra tính hợp lệ của tin nhắn và lưu thông tin bookmark. \\ \hline
Neon Database & Database & Lưu trữ các bản ghi bookmark liên kết giữa User, Message và các thông tin thư mục, ghi chú. \\ \hline
\end{tabular}
\caption{Các thành phần tham gia luồng Đánh dấu cuộc trò chuyện}
\end{table}

\subparagraph{Mô tả chi tiết}
\begin{itemize}
\item Khi đang trò chuyện, User nhấn nút "Đánh dấu" tại một tin nhắn mong muốn. Frontend mở giao diện pop-up để User lựa chọn thư mục lưu trữ và điền ghi chú cá nhân.
\item Frontend gửi yêu cầu `POST /bookmarks` kèm theo `messageId`, `folderName`, và `notes` qua Kong Gateway tới Conversation Service.
\item Conversation Service kiểm tra bản ghi tin nhắn và ghi nhận thông tin bookmark mới vào Neon Database.
\item Sau khi cơ sở dữ liệu ghi nhận thành công, dịch vụ trả về mã phản hồi `HTTP 201 Created` để Frontend cập nhật trạng thái "Đã lưu" trên màn hình của người dùng.
\end{itemize}
